% Original analytic teaching diagram, 2026-11-09.
\definecolor{VB}{HTML}{7188A3}\definecolor{VBF}{HTML}{E5EAF0}
\definecolor{VG}{HTML}{829C8B}\definecolor{VGF}{HTML}{E8EEE9}
\definecolor{VY}{HTML}{B4A16F}\definecolor{VYF}{HTML}{F1ECDD}
\definecolor{VR}{HTML}{A7787D}\definecolor{VRF}{HTML}{F0E5E6}
\definecolor{VP}{HTML}{9588A5}\definecolor{VPF}{HTML}{ECE8F0}
\definecolor{VI}{HTML}{3B4252}\definecolor{VM}{HTML}{6B7280}\definecolor{VL}{HTML}{4C4D4F}
\begin{tikzpicture}[x=1mm,y=1mm,>=Stealth,font=\figurefont\fontsize{8}{10}\selectfont,lab/.style={text=VM,font=\figurefont\fontsize{7.5}{9}\selectfont},box/.style={draw=VB,fill=VBF,rounded corners=1mm,minimum height=11mm,text width=25mm,align=center,text=VI},arr/.style={->,draw=VL,line width=.7pt}]
\node[box] (p) at (0,0) {分布 $\mathbb{P}$\\ 中的随机变量 $\bm{\xi}$};
\node[box,draw=VG,fill=VGF] (s) at (43,0) {固定的 $N$ 项\\ $f_N(\bm{x})$};
\node[box,draw=VY,fill=VYF] (r) at (86,0) {本步读数\\ $\hat{\bm{g}}^{(t)}$};
\draw[arr] (p) -- node[above,lab]{抽样本} node[below,lab]{一次} (s);
\draw[arr] (s) -- node[above,lab]{抽下标} node[below,lab]{每步} (r);
\draw[arr] (p.south) -- ++(0,-7) -| node[pos=.25,below,lab]{每步抽新样本，不经过固定样本} (r.south);
\draw[decorate,decoration={brace,amplitude=2mm},draw=VM,line width=.6pt] (-12,9.5) -- node[above=2.3mm,lab]{数据的随机性：决定解的是哪个问题} (52,9.5);
\draw[decorate,decoration={brace,amplitude=2mm},draw=VM,line width=.6pt] (56,9.5) -- node[above=2.3mm,lab]{求解的随机性} (98,9.5);
\end{tikzpicture}
